\chapter{Symmetry Groups}
\label{chap:symmetry-groups}

% Closed question: What remains invariant under admissible relabeling or
% recompilation?

\begin{chapteressay}

The previous chapter introduced curvature: the algebraic record of
transport's failure to commute. This chapter takes up the dual
question: what does commute? What is invariant under the
relabelings that the curvature exposes as non-trivial?

Invariance is the chapter's subject. Some quantities depend on the
specific labeling of the underlying structure: the coordinates of
a triangle's vertices depend on the choice of coordinate system.
Other quantities do not: the triangle's area is independent of the
coordinate system. The area is invariant under the group of
rigid motions of the plane; the coordinates are not.

The chapter develops invariance through five mathematical
structures.

The first is the basic notion of invariance under relabeling. A
quantity is invariant when its value does not change under any
element of the symmetry group. Different invariants correspond
to different groups: rotations preserve the lengths of vectors;
conformal transformations preserve angles; affine transformations
preserve ratios of distances along lines.

The second is the group action: the mathematical formalization of
what relabelings are admissible. The symmetry group of a square is
the dihedral group $D_4$; the group acts on the square by
permuting vertices. The orbits of the action are the equivalence
classes; the invariants are the functions constant on orbits.

The third is Noether's theorem. For continuous symmetries (Lie
groups) acting on a physical system (with an action functional),
every continuous symmetry corresponds to a conserved quantity.
Time-translation invariance gives energy conservation;
space-translation gives momentum; rotation gives angular
momentum; gauge invariance gives electric charge.

The fourth is the Aharonov-Bohm effect: a topological invariant
manifesting as a measurable phase shift. The magnetic flux
through a closed loop is a topological invariant; the charged
particle's wave function responds to the flux even when the
local field is zero. The effect demonstrates that topology can
matter independently of local geometry.

The fifth is bootstrapping as the algorithmic fixed point. The
Lean 4 compiler is written in Lean 4; recompiling the source
with the compiler produces an identical compiler. The semantic
content is the invariant under recompilation. The bootstrap is
the chapter's algorithmic instance of invariance.

The chapter's ground example is the triangle's area: invariant
under rigid motions. The chapter's second concrete example is
the family shortbread recipe: invariant under unit conversions,
temperature scale changes, and ingredient substitutions. The
recipe survives transport across units, temperatures, brands, and
kitchens; the cookies are recognizably the same regardless of the
specific admissible substitutions used.

The chapter's closed question is:

\begin{center}
\emph{What remains invariant under admissible relabeling or
recompilation?}
\end{center}

The chapter's answer is: the invariants of the group action. The
group specifies the admissible relabelings; the invariants are
the quantities that survive every element of the group. For
geometric quantities, the invariants are the area, the length,
the angle, the curvature scalar, the topological invariants. For
physical quantities, the invariants are the Noether charges and
the topological invariants. For algorithmic quantities, the
invariants are the semantic content of the compiled program.

The chapter's discipline: state claims in terms of invariants.
Claims that depend on specific labelings are incomplete; they
must be reformulated in terms of group-invariant quantities to
be admissible structural statements.

The next chapter is the final chapter. It addresses why the record
only grows --- the append-only property of the mathematical
ledger --- and when the record is closed enough to support
inference. The book closes with the final type-check on the
accumulated formal apparatus.

\end{chapteressay}

\section{Invariance Under Relabeling}
\label{se:invariance-under-relabeling}

A triangle has area $A$. Rotate the triangle by any angle about any
point; reflect it across any line; translate it by any vector. After
each transformation, the triangle's area is still $A$. The area is
invariant under the group of rigid motions of the plane.

The triangle's specific coordinates change with each transformation:
a triangle with vertices $(0, 0), (1, 0), (0, 1)$ becomes, after a
$90^\circ$ rotation, the triangle with vertices $(0, 0), (0, 1),
(-1, 0)$. The vertex coordinates have changed; the relations among
the vertices (specifically, the area) have not.

This is the chapter's first claim. An invariant is a quantity that
does not change under the action of a specified group of
transformations. The area of a planar region is invariant under the
group of rigid motions. The length of a vector is invariant under
rotations. The trace of a matrix is invariant under conjugation. In
each case, the quantity is a function of the underlying structure,
not of the specific labeling.

The invariant is the structural fact; the specific labeling is the
representation. The same structural fact (the triangle's area) can
be expressed in many different coordinate systems; the invariant is
what they all agree on.

The chapter's discipline: invariants are admissible. Specific
labelings are not. A claim that depends on the specific labeling is
not yet admissible as a structural claim about the underlying
object; the claim must be expressed in terms of invariants.

In mathematics, this discipline is pervasive. A vector space has
basis-independent properties (dimension, rank, trace, determinant);
these are admissible structural claims. A specific basis is one of
many possible labelings; claims tied to a specific basis are
incomplete without additional context.

In Lean, the discipline is enforced by the type system. A theorem
about a vector space is stated in basis-independent terms; the
proof may use a specific basis but must produce a result that
holds regardless. The compiler does not directly enforce
invariance, but the typing discipline encourages it.

\section{Group Action}
\label{se:group-action}

The symmetry group of a square consists of eight elements: four
rotations ($0^\circ$, $90^\circ$, $180^\circ$, $270^\circ$) and
four reflections (across the two diagonals and across the two
midlines). This group is called $D_4$, the dihedral group of
order eight.

Each element of $D_4$ is a bijection from the square to itself
that preserves the square's structure: each element maps vertices
to vertices, edges to edges, and the center to the center. The
group is closed under composition: applying two symmetries in
succession gives another symmetry. The identity is the
$0^\circ$ rotation. Each symmetry has an inverse (the reverse
rotation or the same reflection).

A group $G$ acts on a set $X$ when there is a function $G \times X
\to X$ satisfying:
\begin{itemize}
\item Identity: the identity element of $G$ acts as the identity
function on $X$.
\item Composition: for any $g, h \in G$ and $x \in X$, $(g h) \cdot
x = g \cdot (h \cdot x)$.
\end{itemize}
The group action makes precise the idea that the group's elements
relabel the elements of $X$ in a structured way.

For the square, the set $X$ is the square's set of vertices; $D_4$
acts on $X$ by permuting the vertices. The action is faithful
(different group elements give different permutations) and
transitive (any vertex can be mapped to any other).

The chapter's second claim is that the group action is the
mathematical formalization of admissible relabeling. The group
specifies what relabelings are admissible; the action specifies
how each relabeling acts on the underlying set. The invariants
of the action are the quantities preserved by every group
element.

The orbit of a point $x \in X$ under $G$ is the set $\{g \cdot x :
g \in G\}$: all the points reachable from $x$ by some group
element. The orbits partition $X$. Two points in the same orbit
are equivalent under the group's relabeling; they differ only by
a relabeling that the group permits.

The stabilizer of $x$ is the subgroup $\{g \in G : g \cdot x = x\}$:
all elements that fix $x$. The orbit-stabilizer theorem says: the
size of the orbit times the size of the stabilizer equals the size
of the group. For finite groups acting on finite sets, this gives
a useful relation.

In Lean, the group-action structure is the typeclass
\texttt{MulAction} in \texttt{Mathlib.GroupTheory.GroupAction.Basic}:

\begin{leanexcerpt}{MulAction}
class MulAction ($\alpha$ : Type u) ($\beta$ : Type v) [Monoid $\alpha$]
    extends SMul $\alpha$ $\beta$ where
  one_smul : $\forall$ b : $\beta$, (1 : $\alpha$) $\bullet$ b = b
  mul_smul : $\forall$ (a$_1$ a$_2$ : $\alpha$) (b : $\beta$),
    (a$_1$ * a$_2$) $\bullet$ b = a$_1$ $\bullet$ a$_2$ $\bullet$ b
\end{leanexcerpt}

The class formalizes the (monoid) action with the identity and
composition axioms. The compiler enforces both: an instance must
provide proofs of both properties. (Group actions arise as the
special case where the monoid is a group; \texttt{MulAction} is
the general form.)

The chapter's claim is that the group-action structure is the
foundation of symmetry-based analysis. Whenever a system has
admissible relabelings, the relabelings form a group, the group
acts on the system, and the invariants of the action are the
structural facts.

\section{Noether's Theorem}
\label{se:noethers-theorem}

Emmy Noether proved in 1918 that every continuous symmetry of the
action functional corresponds to a conserved quantity. The result
is one of the most important theorems of mathematical physics.

For a mechanical system with Lagrangian $L(q, \dot{q}, t)$ and
action $S[q] = \int L \, dt$, suppose the action is invariant under
a continuous transformation $q \to q + \epsilon \delta q$, where
$\delta q$ is a generator of the transformation and $\epsilon$ is
the parameter. Then the quantity
\[
Q = \frac{\partial L}{\partial \dot{q}} \delta q
\]
is conserved: $dQ/dt = 0$ along solutions of the Euler-Lagrange
equations.

Specific applications:
\begin{itemize}
\item Time-translation invariance ($t \to t + \epsilon$): energy is
conserved.
\item Space-translation invariance ($q \to q + \epsilon$): momentum
is conserved.
\item Rotation invariance: angular momentum is conserved.
\item Gauge invariance: electric charge is conserved.
\end{itemize}

Each symmetry generates a specific conserved quantity. The
correspondence is structural: the symmetry's generator (the
$\delta q$) becomes the conjugate quantity to the canonical momentum
$\partial L / \partial \dot{q}$.

The chapter's third claim is that conserved quantities are
invariants of the action. The action's symmetries are the
admissible relabelings; the conserved quantities are the
structural invariants. Noether's theorem is the precise
mathematical correspondence between the two.

In Lean, Noether's theorem requires the variational calculus
infrastructure of Chapter 6 plus the group action infrastructure
of this chapter. The combination is non-trivial; \texttt{Mathlib}
has partial support. The full Noether theorem at the level of
field theories is beyond the current state of the Lean
formalization; it is in the chapter's
\texttt{future\_work} category.

\begin{phenom}{Noether's Theorem}
\label{ph:noethers-theorem}

\PhStatement For every continuous symmetry of the action functional
of a physical system, there is a corresponding conserved quantity.

\PhOrigin Emmy Noether proved this in her 1918 paper ``Invariante
Variationsprobleme,'' addressing concerns about energy conservation
in general relativity.

\PhObservation The pendulum's motion conserves energy because its
Lagrangian is time-independent. A spinning top with rotational
symmetry conserves angular momentum about the symmetry axis. A
charged particle in an electromagnetic field has conserved electric
charge due to gauge invariance.

\PhConstraint The theorem requires the symmetry to be continuous
(a Lie group, not a discrete group). Discrete symmetries
correspond to selection rules but not to conserved Noether
charges.

\PhConsequence Conservation laws are not contingent facts; they
follow from the action's symmetry structure. Without the
symmetry, the conservation fails; without the conservation, the
symmetry is not a symmetry of the action.

\end{phenom}

The Noether's Theorem phenomenon is the chapter's load-bearing
structural result. Symmetries and conservations are related by
the theorem; one implies the other.

\section{Aharonov-Bohm}
\label{se:aharonov-bohm}

The Aharonov-Bohm effect, predicted in 1959 and experimentally
verified in the 1980s, shows that a charged particle can be
affected by an electromagnetic field even when the particle passes
through a region where the field is zero.

The setup: a long solenoid carries an electric current, producing
a magnetic field $\mathbf{B}$ inside the solenoid and zero magnetic
field outside. The vector potential $\mathbf{A}$ associated with
the field, however, is nonzero outside the solenoid: the
circulation $\oint \mathbf{A} \cdot d\mathbf{l}$ around any closed
loop enclosing the solenoid equals the magnetic flux $\Phi$ through
the solenoid.

A charged particle traveling on a path that encloses the solenoid
acquires a quantum-mechanical phase shift proportional to $\Phi$.
The phase shift is measurable in an interference experiment: two
beams of electrons, one passing on each side of the solenoid, are
recombined; the interference pattern shifts by an amount
corresponding to the enclosed flux.

The chapter's fourth claim is that the Aharonov-Bohm effect is the
chapter's instance of a topological invariant. The phase shift
depends on the enclosed magnetic flux, not on the local
electromagnetic field at the particle's location. The flux is a
topological quantity (an integral over the surface bounded by the
loop). The local field is the geometric quantity. The particle
responds to the topology, not the local geometry.

The mathematical formulation uses gauge theory. The
electromagnetic field is described by a connection on a $U(1)$
bundle over spacetime. The connection's curvature is the
electromagnetic field tensor. The bundle's topology is captured
by the integer-valued Chern number, which for a magnetic flux
$\Phi$ on a closed surface gives $\Phi / (h/e)$ (where $h/e$ is
the magnetic flux quantum).

The chapter's claim is that the Aharonov-Bohm effect demonstrates
the existence of topologically invariant quantities that are not
captured by the local geometric data. The enclosed flux is a
topological invariant; the local field at the particle is the
geometric quantity. The particle's phase shift is the
chapter's algebraic record of the topological invariant.

This is a deep claim. It says: some physical effects depend on
global topological structure, not on local geometric structure.
The chapter's structures (invariants under group action,
Noether charges, topological invariants) include such global
quantities; the chapter's discipline is to recognize them as
admissible structural facts.

In Lean, the analog appears in the typeclass machinery for
topological invariants. The chapter's
\texttt{INVARIANT} typeclass (informally introduced via the
preserved invariants of the group action) would carry the
admissibility of topological invariants. The full development
is beyond the chapter's current Lean infrastructure; the
chapter notes this as a future development.

\section{Bootstrapping as Fixed Point}
\label{se:bootstrapping-as-fixed-point}

The Lean 4 compiler is written in Lean 4. The first version of the
compiler was written in another language (C++), but once a working
Lean 4 compiler existed, the source could be compiled by the
compiler to produce a new compiler binary. The new binary should
be identical to the original.

This is the bootstrapping property. A compiler is a fixed point of
its own compilation process: $\text{stage}(N+1) = \text{stage}(N)$.
The fixed point is the invariant under recompilation: the
compiler does not change when recompiled with itself.

The fixed point is the chapter's algorithmic instance of invariance.
The compiler's binary changes its internal representation when
recompiled (different memory layout, different timestamps), but
the semantic content is invariant: the same program compiles to a
program that behaves the same way. The semantic content is the
invariant under recompilation.

For Lean 4, the bootstrapping was achieved in 2021. The compiler's
source compiles itself to a binary that, when used to compile the
source again, produces a semantically equivalent binary --- a
fixed point up to the chosen build equivalence. (Bit-for-bit
identity is not guaranteed across all build environments;
timestamps, layout, and metadata vary. The invariant is semantic
behavior, not byte sequence.) The verification is the chapter's
algorithmic version of the bootstrap fixed point.

The chapter's fifth claim is that bootstrapping is the
algorithmic form of invariance. The compiler is the algorithm; the
recompilation is the group action (recompile by the same compiler);
the bootstrap fixed point is the invariant element.

In Lean, this is implicit in the language's design. The
\texttt{COMPILED} typeclass certifies that a process has reached
its compiled form. The
\texttt{ElaborationProcess} typeclass tracks the elaboration
through load $\to$ transform $\to$ boolean stages; the boolean
stage is \texttt{HALTED}. The bootstrap is the demonstration that
the compiler's source plus a working compiler produces the
compiler again.

\begin{leanexcerpt}{ComputerProgram}
inductive ComputerProgram
  | load      : Fact $\to$ Prop $\to$ Type $\to$ ComputerProgram
  | transform : Fact $\to$ Fact $\to$ Prop $\to$ Prop
              $\to$ Type $\to$ Type 1
              $\to$ ComputerProgram $\to$ ComputerProgram
  | boolean   : Fact $\to$ Fact $\to$ Fact $\to$ Prop $\to$ Prop $\to$ Prop
              $\to$ Type $\to$ Type 1 $\to$ Type i
              $\to$ ComputerProgram $\to$ ComputerProgram
\end{leanexcerpt}

The project's \texttt{ComputerProgram} inductive in
\texttt{ComputerProgram.lean} captures the bootstrap layering: a
program is built from \texttt{load} (the initial Fact/Prop/Type
triple), \texttt{transform} (an admissible rewrite step that
extends a prior program), and \texttt{boolean} (a step that
decides between admissible branches). The recompilation is the
fixed point inside this inductive: a closed compilation
\texttt{ComputerProgram} that, when fed back to itself, produces
the same shape under the project's heartbeat calibration. The
compiler is the invariant under its own action.

The chapter closes here. Invariance under relabeling is the
chapter's first structure: quantities that do not change under
admissible relabelings. Group action formalizes the relabelings.
Noether's theorem connects continuous symmetries to conserved
quantities. The Aharonov-Bohm effect demonstrates topological
invariants. Bootstrapping is the algorithmic fixed point: the
compiler is invariant under its own recompilation.

\begin{bridge}

The chapter's question was what remains invariant under admissible
relabeling or recompilation.

The answer: the invariants of the group action. The group of
admissible relabelings acts on the underlying record; the
invariants of the action are the structural facts. The area
of a triangle is invariant under rigid motions; the trace of a
matrix is invariant under conjugation; the conserved Noether
charges are invariant under continuous symmetries; the
topological invariants (Aharonov-Bohm flux quantum, Chern numbers)
are invariant under continuous deformations; the compiler's
semantic content is invariant under recompilation.

What remains is closure. The chapter has identified what does not
change; the next chapter takes up what grows. Chapter 11 is the
final chapter: it addresses why the record only grows and when
the record is closed enough to support inference.

\end{bridge}

\begin{coda}{The Recipe Surviving Substitution}

The recipe is for shortbread cookies. It comes from a family
cookbook compiled in the 1950s by a great-grandmother who passed
the cookbook to her daughter, who passed it to her daughter, who
passed it to the current cook. The recipe has been used dozens of
times across four generations.

The recipe is written in imperial units: 1 cup of butter, 1/2 cup
of sugar, 2 cups of flour, 1/4 teaspoon of salt. The instructions
are spare: cream the butter and sugar, mix in flour and salt,
press into a pan, bake at 325 degrees Fahrenheit for 45 minutes,
cool before cutting.

The current cook is in Paris, attending a year-long visiting
position. The kitchen is small, the oven is in Celsius, the
units on the packages are grams, the flour available is not quite
the same flour the family has always used. The cook adapts.

Cup to gram: 1 cup of butter is approximately 227 grams; 1/2 cup of
sugar is approximately 100 grams; 2 cups of flour is approximately
250 grams; 1/4 teaspoon of salt is approximately 1.5 grams.
Fahrenheit to Celsius: 325 F is approximately 163 C. The cook
makes the substitutions and prepares the dough.

The flour available is French type 55, similar but not identical
to American all-purpose flour. The cook uses it. The butter is
European unsalted butter, with slightly higher fat content than
American butter. The cook uses it. The salt is fine sea salt
rather than table salt; the cook accounts for the difference in
crystal size by using slightly less.

The shortbread bakes for 45 minutes. The cook takes it out, cools
it, cuts it into squares. The first square is taken to taste.

The taste is recognizable. The texture is familiar. The shortbread
is the same shortbread the cook grew up with. Slightly different
in subtle ways --- the European butter contributes a slightly
different mouthfeel; the French flour gives a slightly different
crumb --- but recognizably the same recipe, the same family
tradition.

What survived the substitution? Not the specific brand of butter,
not the American cup measure, not the Fahrenheit temperature, not
the American flour. What survived is the structural relationship
of the ingredients: the ratio of butter to sugar to flour, the
short development of gluten, the slow even bake, the
post-baking cool. The recipe is the invariant under the
unit conversions and brand substitutions.

This is the chapter's invariance in domestic form. The recipe has
been transported across units (cup to gram), across temperatures
(Fahrenheit to Celsius), across brands (American flour to French
flour), and across geographies (American kitchen to Parisian
kitchen). The invariant under all of these admissible
substitutions is the shortbread itself: the taste, the texture,
the result.

The chapter's group action is the substitution group: rescalings
of units, conversions of temperatures, substitutions of equivalent
ingredients. The recipe's invariant is the shortbread. The
shortbread is recognized as the same regardless of which specific
admissible substitution path was taken to produce it.

If the substitutions are inadmissible --- a different temperature,
a different ingredient that does not bake the same way, a
different ratio of components --- the recipe fails. The cookie is
not the shortbread; the substitution has left the equivalence
class. The chapter's discipline appears: admissible substitutions
preserve the invariant; inadmissible ones do not.

The Lean 4 bootstrapping example from the chapter is the
algorithmic analog. The compiler is the recipe. Recompilation
is the substitution: the compiler's source is recompiled with
itself; the new binary is the bootstrap fixed point. The
compiler is invariant under its own recompilation.

The grandmother's shortbread is the chapter's symmetry group at
the level of family cooking. The recipe has been passed down four
generations; each generation has adapted it slightly to the
available ingredients and equipment; the invariant is the
shortbread itself. The recipe is the chapter's recipe-as-invariant:
the structural fact preserved through admissible substitutions.

When the cook returns to the United States in a year, the recipe
will be made again, this time in the original American kitchen
with the original American ingredients. The cookies will be the
same cookies. The recipe will be the same recipe. The
bootstrapping fixed point will hold.

\end{coda}
